\chapter*{Resumen}
\addcontentsline{toc}{chapter}{Resumen}
\markboth{Resumen}{Resumen}

\begin{spacing}{1.0}

Este Trabajo de Fin de Grado estudia el análisis automático de electrocardiogramas disponibles como imágenes, no como series temporales originales. El caso de estudio es la sospecha de síndrome de Brugada, una enfermedad rara para la que resulta difícil reunir conjuntos amplios de pacientes etiquetados. El objetivo es comparar distintas formas de aprovechar pocos ejemplos para priorizar ECG que requieren revisión especializada.

Se comparan dos enfoques: una red convolucional entrenada de forma supervisada y un modelo de visión y lenguaje con aprendizaje en contexto. Ambos reciben una imagen V1 y se evalúan por paciente para \(K=\{2,4,8,16,32\}\). Los experimentos utilizan un simulador QRS/ST con 100 pacientes y una cohorte procesada del Hospital Universitario Central de Asturias (HUCA) con 317.

La ResNet18 aprende tres hallazgos morfológicos a partir de pacientes del simulador. Su evaluación sobre el HUCA mide transferencia sintético-real, ya que la cohorte no contiene anotaciones independientes de RBBB, elevación del ST e inversión de T para entrenar esa cabeza multietiqueta.

Gemma~4 utiliza los ejemplos como demostraciones en el prompt, sin modificar sus pesos. Se comparan zero-shot, ICL estándar con cobertura mínima de ambas clases e ICL balanceado, utilizando contexto sintético o clínico.

En el simulador, la CNN alcanza una \textit{Balanced Accuracy} (BA) de \(0,885 \pm 0,013\) en \(K=32\); el ICL balanceado llega a \(0,904 \pm 0,011\), frente a 0,500 en zero-shot. Sobre el HUCA, la CNN obtiene su máximo, \(0,551 \pm 0,034\), en \(K=4\). MMD eleva el punto de \(K=32\) a \(0,546 \pm 0,011\), sin superar ese máximo. Gemma~4 con demostraciones sintéticas permanece cerca de una predicción negativa constante y alcanza \(0,514 \pm 0,008\) en \(K=16\). Con demostraciones de la propia cohorte, el ICL balanceado llega a \(0,761 \pm 0,011\) en \(K=32\), con sensibilidad de 0,741 y especificidad de 0,781.

Los dos enfoques aprenden la regla del dominio sintético, pero su transferencia mediante ejemplos del simulador es limitada. El aprendizaje en contexto obtiene su mejor resultado cuando las demostraciones representan la tarea clínica local. Ningún modelo dispone aún de evidencia suficiente para un triaje fiable. La principal aportación consiste en cuantificar esta brecha mediante un protocolo documentado y señalar la necesidad de anotaciones morfológicas clínicas, más casos positivos y validación externa.

\textbf{Palabras clave:} electrocardiograma como imagen, síndrome de Brugada, pocos datos etiquetados, aprendizaje supervisado, redes neuronales convolucionales, modelos de visión y lenguaje, aprendizaje en contexto, triaje clínico.

\end{spacing}
